% Fallback bibliography: the same 30 corrected references as an
% inline thebibliography environment, for compiling without BibTeX.
% To use, replace the \bibliographystyle/\bibliography lines in
% main.tex with:  % Fallback bibliography: the same 30 corrected references as an
% inline thebibliography environment, for compiling without BibTeX.
% To use, replace the \bibliographystyle/\bibliography lines in
% main.tex with:  % Fallback bibliography: the same 30 corrected references as an
% inline thebibliography environment, for compiling without BibTeX.
% To use, replace the \bibliographystyle/\bibliography lines in
% main.tex with:  % Fallback bibliography: the same 30 corrected references as an
% inline thebibliography environment, for compiling without BibTeX.
% To use, replace the \bibliographystyle/\bibliography lines in
% main.tex with:  \input{references-inline}

\begin{thebibliography}{30}

\bibitem{nguyen2024}
H.~T. Nguyen, M.~Usman, R.~Buyya,
iQuantum: A toolkit for modeling and simulation of quantum computing
environments,
Software: Practice and Experience 54~(6) (2024) 1141--1171.

\bibitem{calheiros2011}
R.~N. Calheiros, R.~Ranjan, A.~Beloglazov, C.~A.~F. De~Rose, R.~Buyya,
CloudSim: A toolkit for modeling and simulation of cloud computing environments
and evaluation of resource provisioning algorithms,
Software: Practice and Experience 41~(1) (2011) 23--50.

\bibitem{mahmud2022}
R.~Mahmud, S.~Pallewatta, M.~Goudarzi, R.~Buyya,
iFogSim2: An extended iFogSim simulator for mobility, clustering, and
microservice management in edge and fog computing environments,
Journal of Systems and Software 190 (2022) 111351.

\bibitem{preskill2018}
J.~Preskill,
Quantum computing in the NISQ era and beyond,
Quantum 2 (2018) 79.

\bibitem{nielsen2010}
M.~A. Nielsen, I.~L. Chuang,
Quantum Computation and Quantum Information,
Cambridge University Press, Cambridge, UK, 2010.

\bibitem{cross2019}
A.~W. Cross, L.~S. Bishop, S.~Sheldon, P.~D. Nation, J.~M. Gambetta,
Validating quantum computers using randomized model circuits,
Physical Review A 100~(3) (2019).

\bibitem{proctor2025}
T.~Proctor, K.~Young, A.~D. Baczewski, R.~Blume-Kohout,
Benchmarking quantum computers,
Nature Reviews Physics 7 (2025) 105--118.

\bibitem{magesan2011}
E.~Magesan, J.~M. Gambetta, J.~Emerson,
Scalable and robust randomized benchmarking of quantum processes,
Physical Review Letters 106~(18) (2011).

\bibitem{magesan2012}
E.~Magesan, J.~M. Gambetta, B.~R. Johnson, C.~A. Ryan, J.~M. Chow,
S.~T. Merkel, M.~P. da~Silva, G.~A. Keefe, M.~B. Rothwell, T.~A. Ohki,
M.~B. Ketchen, J.~M. Steffen,
Efficient measurement of quantum gate error by interleaved randomized
benchmarking,
Physical Review Letters 109~(8) (2012).

\bibitem{yao2025}
G.~Yao, L.~Gupta,
A benchmarking framework for hybrid quantum--classical edge-cloud computing
systems,
Applied Sciences 15 (2025) 10245.

\bibitem{furutanpey2023}
A.~Furutanpey, J.~Barzen, M.~Bechtold, S.~Dustdar, F.~Leymann, P.~Raith,
F.~Truger,
Architectural vision for quantum computing in the edge-cloud continuum,
in: Proc. IEEE International Conference on Quantum Software (QSW), 2023,
pp. 88--103.

\bibitem{mastroianni2023}
C.~Mastroianni, L.~Scarcello, A.~Vinci,
Quantum computing management of a cloud/edge architecture,
in: Proc. ACM International Conference on Computing Frontiers (CF), 2023,
pp. 193--196.

\bibitem{xu2025}
M.~Xu, D.~Niyato, J.~Kang, Z.~Xiong, M.~Chen, D.~I. Kim, X.~Shen,
Hybrid reinforcement learning-based sustainable multi-user computation
offloading for mobile edge-quantum computing,
arXiv:2504.08134 (2025).

\bibitem{huang2024}
S.~Huang, L.~Wang, X.~Wang, B.~Tan, W.~Ni, K.-K. Wong,
Edge intelligence in satellite-terrestrial networks with hybrid quantum
computing,
arXiv:2409.19869 (2024).

\bibitem{satya2017}
M.~Satyanarayanan,
The emergence of edge computing,
Computer 50~(1) (2017) 30--39.

\bibitem{shi2016}
W.~Shi, J.~Cao, Q.~Zhang, Y.~Li, L.~Xu,
Edge computing: Vision and challenges,
IEEE Internet of Things Journal 3~(5) (2016) 637--646.

\bibitem{mao2017}
Y.~Mao, C.~You, J.~Zhang, K.~Huang, K.~B. Letaief,
A survey on mobile edge computing: The communication perspective,
IEEE Communications Surveys \& Tutorials 19~(4) (2017) 2322--2358.

\bibitem{abbas2018}
N.~Abbas, Y.~Zhang, A.~Taherkordi, T.~Skeie,
Mobile edge computing: A survey,
IEEE Internet of Things Journal 5~(1) (2018) 450--465.

\bibitem{leymann2020}
F.~Leymann, J.~Barzen, M.~Falkenthal, D.~Vietz, B.~Weder, K.~Wild,
Quantum in the cloud: Application potentials and research opportunities,
in: Proc. 10th International Conference on Cloud Computing and Services Science
(CLOSER), SciTePress, 2020, pp. 9--24.

\bibitem{weder2021a}
B.~Weder, J.~Barzen, F.~Leymann, M.~Salm, D.~Vietz,
The quantum software lifecycle,
in: Proc. 1st ACM SIGSOFT International Workshop on Architectures and Paradigms
for Engineering Quantum Software (APEQS), ACM, 2020, pp. 2--9.

\bibitem{weder2021b}
% REPLACED: the entry previously here ("Quantum application archives:
% Supporting the reuse of quantum applications") could not be located in any
% bibliographic database and appears not to exist. Substituted with a real,
% topically adjacent paper by the same group -- VERIFY this is what you meant
% to cite, or delete this entry (it is currently uncited in the text).
B.~Weder, J.~Barzen, F.~Leymann, M.~Zimmermann,
Hybrid quantum applications need two orchestrations in superposition: A software
architecture perspective,
in: Proc. 18th IEEE International Conference on Web Services (ICWS), IEEE, 2021.

\bibitem{wild2021}
K.~Wild, U.~Breitenb\"{u}cher, L.~Harzenetter, F.~Leymann, D.~Vietz,
M.~Zimmermann,
TOSCA4QC: Two modeling styles for TOSCA to automate the deployment and
orchestration of quantum applications,
in: Proc. 24th IEEE International Enterprise Distributed Object Computing
Conference (EDOC), IEEE, 2020, pp. 125--134.

\bibitem{garcia2022}
J.~Garcia-Alonso, J.~Rojo, D.~Valencia, E.~Moguel, J.~Berrocal, J.~M. Murillo,
Quantum software as a service through a quantum API gateway,
IEEE Internet Computing 26~(1) (2022) 34--41.

\bibitem{beisel2023}
M.~Beisel, J.~Barzen, S.~Garhofer, F.~Leymann, F.~Truger, B.~Weder,
V.~Yussupov,
Quokka: A service ecosystem for workflow-based execution of variational quantum
algorithms,
in: Service-Oriented Computing -- ICSOC 2022 Workshops, Vol. 13821 of LNCS,
Springer, 2023, pp. 369--373.

\bibitem{quetschlich2023}
N.~Quetschlich, L.~Burgholzer, R.~Wille,
MQT Bench: Benchmarking software and design automation tools for quantum
computing,
Quantum 7 (2023) 1062.

\bibitem{li2020}
A.~Li, S.~Stein, S.~Krishnamoorthy, J.~Ang,
QASMBench: A low-level quantum benchmark suite for NISQ evaluation and
simulation,
ACM Transactions on Quantum Computing 4~(2) (2023) 1--26.

\bibitem{tomesh2022}
T.~Tomesh, P.~Gokhale, V.~Omole, G.~S. Ravi, K.~N. Smith, J.~Viszlai,
X.-C. Wu, N.~Hardavellas, M.~R. Martonosi, F.~T. Chong,
SupermarQ: A scalable quantum benchmark suite,
in: Proc. IEEE International Symposium on High-Performance Computer
Architecture (HPCA), IEEE, 2022, pp. 587--603.

\bibitem{finzgar2024}
J.~R. Fin\v{z}gar, P.~Ross, L.~H\"{o}lscher, J.~Klepsch, A.~Luckow,
QUARK: A framework for quantum computing application benchmarking,
in: Proc. IEEE International Conference on Quantum Computing and Engineering
(QCE), IEEE, 2022, pp. 226--237.

\bibitem{farhi2014}
E.~Farhi, J.~Goldstone, S.~Gutmann,
A quantum approximate optimization algorithm,
arXiv:1411.4028 (2014).

\bibitem{cerezo2021}
M.~Cerezo, A.~Arrasmith, R.~Babbush, S.~C. Benjamin, S.~Endo, K.~Fujii,
J.~R. McClean, K.~Mitarai, X.~Yuan, L.~Cincio, P.~J. Coles,
Variational quantum algorithms,
Nature Reviews Physics 3 (2021) 625--644.

\end{thebibliography}


\begin{thebibliography}{30}

\bibitem{nguyen2024}
H.~T. Nguyen, M.~Usman, R.~Buyya,
iQuantum: A toolkit for modeling and simulation of quantum computing
environments,
Software: Practice and Experience 54~(6) (2024) 1141--1171.

\bibitem{calheiros2011}
R.~N. Calheiros, R.~Ranjan, A.~Beloglazov, C.~A.~F. De~Rose, R.~Buyya,
CloudSim: A toolkit for modeling and simulation of cloud computing environments
and evaluation of resource provisioning algorithms,
Software: Practice and Experience 41~(1) (2011) 23--50.

\bibitem{mahmud2022}
R.~Mahmud, S.~Pallewatta, M.~Goudarzi, R.~Buyya,
iFogSim2: An extended iFogSim simulator for mobility, clustering, and
microservice management in edge and fog computing environments,
Journal of Systems and Software 190 (2022) 111351.

\bibitem{preskill2018}
J.~Preskill,
Quantum computing in the NISQ era and beyond,
Quantum 2 (2018) 79.

\bibitem{nielsen2010}
M.~A. Nielsen, I.~L. Chuang,
Quantum Computation and Quantum Information,
Cambridge University Press, Cambridge, UK, 2010.

\bibitem{cross2019}
A.~W. Cross, L.~S. Bishop, S.~Sheldon, P.~D. Nation, J.~M. Gambetta,
Validating quantum computers using randomized model circuits,
Physical Review A 100~(3) (2019).

\bibitem{proctor2025}
T.~Proctor, K.~Young, A.~D. Baczewski, R.~Blume-Kohout,
Benchmarking quantum computers,
Nature Reviews Physics 7 (2025) 105--118.

\bibitem{magesan2011}
E.~Magesan, J.~M. Gambetta, J.~Emerson,
Scalable and robust randomized benchmarking of quantum processes,
Physical Review Letters 106~(18) (2011).

\bibitem{magesan2012}
E.~Magesan, J.~M. Gambetta, B.~R. Johnson, C.~A. Ryan, J.~M. Chow,
S.~T. Merkel, M.~P. da~Silva, G.~A. Keefe, M.~B. Rothwell, T.~A. Ohki,
M.~B. Ketchen, J.~M. Steffen,
Efficient measurement of quantum gate error by interleaved randomized
benchmarking,
Physical Review Letters 109~(8) (2012).

\bibitem{yao2025}
G.~Yao, L.~Gupta,
A benchmarking framework for hybrid quantum--classical edge-cloud computing
systems,
Applied Sciences 15 (2025) 10245.

\bibitem{furutanpey2023}
A.~Furutanpey, J.~Barzen, M.~Bechtold, S.~Dustdar, F.~Leymann, P.~Raith,
F.~Truger,
Architectural vision for quantum computing in the edge-cloud continuum,
in: Proc. IEEE International Conference on Quantum Software (QSW), 2023,
pp. 88--103.

\bibitem{mastroianni2023}
C.~Mastroianni, L.~Scarcello, A.~Vinci,
Quantum computing management of a cloud/edge architecture,
in: Proc. ACM International Conference on Computing Frontiers (CF), 2023,
pp. 193--196.

\bibitem{xu2025}
M.~Xu, D.~Niyato, J.~Kang, Z.~Xiong, M.~Chen, D.~I. Kim, X.~Shen,
Hybrid reinforcement learning-based sustainable multi-user computation
offloading for mobile edge-quantum computing,
arXiv:2504.08134 (2025).

\bibitem{huang2024}
S.~Huang, L.~Wang, X.~Wang, B.~Tan, W.~Ni, K.-K. Wong,
Edge intelligence in satellite-terrestrial networks with hybrid quantum
computing,
arXiv:2409.19869 (2024).

\bibitem{satya2017}
M.~Satyanarayanan,
The emergence of edge computing,
Computer 50~(1) (2017) 30--39.

\bibitem{shi2016}
W.~Shi, J.~Cao, Q.~Zhang, Y.~Li, L.~Xu,
Edge computing: Vision and challenges,
IEEE Internet of Things Journal 3~(5) (2016) 637--646.

\bibitem{mao2017}
Y.~Mao, C.~You, J.~Zhang, K.~Huang, K.~B. Letaief,
A survey on mobile edge computing: The communication perspective,
IEEE Communications Surveys \& Tutorials 19~(4) (2017) 2322--2358.

\bibitem{abbas2018}
N.~Abbas, Y.~Zhang, A.~Taherkordi, T.~Skeie,
Mobile edge computing: A survey,
IEEE Internet of Things Journal 5~(1) (2018) 450--465.

\bibitem{leymann2020}
F.~Leymann, J.~Barzen, M.~Falkenthal, D.~Vietz, B.~Weder, K.~Wild,
Quantum in the cloud: Application potentials and research opportunities,
in: Proc. 10th International Conference on Cloud Computing and Services Science
(CLOSER), SciTePress, 2020, pp. 9--24.

\bibitem{weder2021a}
B.~Weder, J.~Barzen, F.~Leymann, M.~Salm, D.~Vietz,
The quantum software lifecycle,
in: Proc. 1st ACM SIGSOFT International Workshop on Architectures and Paradigms
for Engineering Quantum Software (APEQS), ACM, 2020, pp. 2--9.

\bibitem{weder2021b}
% REPLACED: the entry previously here ("Quantum application archives:
% Supporting the reuse of quantum applications") could not be located in any
% bibliographic database and appears not to exist. Substituted with a real,
% topically adjacent paper by the same group -- VERIFY this is what you meant
% to cite, or delete this entry (it is currently uncited in the text).
B.~Weder, J.~Barzen, F.~Leymann, M.~Zimmermann,
Hybrid quantum applications need two orchestrations in superposition: A software
architecture perspective,
in: Proc. 18th IEEE International Conference on Web Services (ICWS), IEEE, 2021.

\bibitem{wild2021}
K.~Wild, U.~Breitenb\"{u}cher, L.~Harzenetter, F.~Leymann, D.~Vietz,
M.~Zimmermann,
TOSCA4QC: Two modeling styles for TOSCA to automate the deployment and
orchestration of quantum applications,
in: Proc. 24th IEEE International Enterprise Distributed Object Computing
Conference (EDOC), IEEE, 2020, pp. 125--134.

\bibitem{garcia2022}
J.~Garcia-Alonso, J.~Rojo, D.~Valencia, E.~Moguel, J.~Berrocal, J.~M. Murillo,
Quantum software as a service through a quantum API gateway,
IEEE Internet Computing 26~(1) (2022) 34--41.

\bibitem{beisel2023}
M.~Beisel, J.~Barzen, S.~Garhofer, F.~Leymann, F.~Truger, B.~Weder,
V.~Yussupov,
Quokka: A service ecosystem for workflow-based execution of variational quantum
algorithms,
in: Service-Oriented Computing -- ICSOC 2022 Workshops, Vol. 13821 of LNCS,
Springer, 2023, pp. 369--373.

\bibitem{quetschlich2023}
N.~Quetschlich, L.~Burgholzer, R.~Wille,
MQT Bench: Benchmarking software and design automation tools for quantum
computing,
Quantum 7 (2023) 1062.

\bibitem{li2020}
A.~Li, S.~Stein, S.~Krishnamoorthy, J.~Ang,
QASMBench: A low-level quantum benchmark suite for NISQ evaluation and
simulation,
ACM Transactions on Quantum Computing 4~(2) (2023) 1--26.

\bibitem{tomesh2022}
T.~Tomesh, P.~Gokhale, V.~Omole, G.~S. Ravi, K.~N. Smith, J.~Viszlai,
X.-C. Wu, N.~Hardavellas, M.~R. Martonosi, F.~T. Chong,
SupermarQ: A scalable quantum benchmark suite,
in: Proc. IEEE International Symposium on High-Performance Computer
Architecture (HPCA), IEEE, 2022, pp. 587--603.

\bibitem{finzgar2024}
J.~R. Fin\v{z}gar, P.~Ross, L.~H\"{o}lscher, J.~Klepsch, A.~Luckow,
QUARK: A framework for quantum computing application benchmarking,
in: Proc. IEEE International Conference on Quantum Computing and Engineering
(QCE), IEEE, 2022, pp. 226--237.

\bibitem{farhi2014}
E.~Farhi, J.~Goldstone, S.~Gutmann,
A quantum approximate optimization algorithm,
arXiv:1411.4028 (2014).

\bibitem{cerezo2021}
M.~Cerezo, A.~Arrasmith, R.~Babbush, S.~C. Benjamin, S.~Endo, K.~Fujii,
J.~R. McClean, K.~Mitarai, X.~Yuan, L.~Cincio, P.~J. Coles,
Variational quantum algorithms,
Nature Reviews Physics 3 (2021) 625--644.

\end{thebibliography}


\begin{thebibliography}{30}

\bibitem{nguyen2024}
H.~T. Nguyen, M.~Usman, R.~Buyya,
iQuantum: A toolkit for modeling and simulation of quantum computing
environments,
Software: Practice and Experience 54~(6) (2024) 1141--1171.

\bibitem{calheiros2011}
R.~N. Calheiros, R.~Ranjan, A.~Beloglazov, C.~A.~F. De~Rose, R.~Buyya,
CloudSim: A toolkit for modeling and simulation of cloud computing environments
and evaluation of resource provisioning algorithms,
Software: Practice and Experience 41~(1) (2011) 23--50.

\bibitem{mahmud2022}
R.~Mahmud, S.~Pallewatta, M.~Goudarzi, R.~Buyya,
iFogSim2: An extended iFogSim simulator for mobility, clustering, and
microservice management in edge and fog computing environments,
Journal of Systems and Software 190 (2022) 111351.

\bibitem{preskill2018}
J.~Preskill,
Quantum computing in the NISQ era and beyond,
Quantum 2 (2018) 79.

\bibitem{nielsen2010}
M.~A. Nielsen, I.~L. Chuang,
Quantum Computation and Quantum Information,
Cambridge University Press, Cambridge, UK, 2010.

\bibitem{cross2019}
A.~W. Cross, L.~S. Bishop, S.~Sheldon, P.~D. Nation, J.~M. Gambetta,
Validating quantum computers using randomized model circuits,
Physical Review A 100~(3) (2019).

\bibitem{proctor2025}
T.~Proctor, K.~Young, A.~D. Baczewski, R.~Blume-Kohout,
Benchmarking quantum computers,
Nature Reviews Physics 7 (2025) 105--118.

\bibitem{magesan2011}
E.~Magesan, J.~M. Gambetta, J.~Emerson,
Scalable and robust randomized benchmarking of quantum processes,
Physical Review Letters 106~(18) (2011).

\bibitem{magesan2012}
E.~Magesan, J.~M. Gambetta, B.~R. Johnson, C.~A. Ryan, J.~M. Chow,
S.~T. Merkel, M.~P. da~Silva, G.~A. Keefe, M.~B. Rothwell, T.~A. Ohki,
M.~B. Ketchen, J.~M. Steffen,
Efficient measurement of quantum gate error by interleaved randomized
benchmarking,
Physical Review Letters 109~(8) (2012).

\bibitem{yao2025}
G.~Yao, L.~Gupta,
A benchmarking framework for hybrid quantum--classical edge-cloud computing
systems,
Applied Sciences 15 (2025) 10245.

\bibitem{furutanpey2023}
A.~Furutanpey, J.~Barzen, M.~Bechtold, S.~Dustdar, F.~Leymann, P.~Raith,
F.~Truger,
Architectural vision for quantum computing in the edge-cloud continuum,
in: Proc. IEEE International Conference on Quantum Software (QSW), 2023,
pp. 88--103.

\bibitem{mastroianni2023}
C.~Mastroianni, L.~Scarcello, A.~Vinci,
Quantum computing management of a cloud/edge architecture,
in: Proc. ACM International Conference on Computing Frontiers (CF), 2023,
pp. 193--196.

\bibitem{xu2025}
M.~Xu, D.~Niyato, J.~Kang, Z.~Xiong, M.~Chen, D.~I. Kim, X.~Shen,
Hybrid reinforcement learning-based sustainable multi-user computation
offloading for mobile edge-quantum computing,
arXiv:2504.08134 (2025).

\bibitem{huang2024}
S.~Huang, L.~Wang, X.~Wang, B.~Tan, W.~Ni, K.-K. Wong,
Edge intelligence in satellite-terrestrial networks with hybrid quantum
computing,
arXiv:2409.19869 (2024).

\bibitem{satya2017}
M.~Satyanarayanan,
The emergence of edge computing,
Computer 50~(1) (2017) 30--39.

\bibitem{shi2016}
W.~Shi, J.~Cao, Q.~Zhang, Y.~Li, L.~Xu,
Edge computing: Vision and challenges,
IEEE Internet of Things Journal 3~(5) (2016) 637--646.

\bibitem{mao2017}
Y.~Mao, C.~You, J.~Zhang, K.~Huang, K.~B. Letaief,
A survey on mobile edge computing: The communication perspective,
IEEE Communications Surveys \& Tutorials 19~(4) (2017) 2322--2358.

\bibitem{abbas2018}
N.~Abbas, Y.~Zhang, A.~Taherkordi, T.~Skeie,
Mobile edge computing: A survey,
IEEE Internet of Things Journal 5~(1) (2018) 450--465.

\bibitem{leymann2020}
F.~Leymann, J.~Barzen, M.~Falkenthal, D.~Vietz, B.~Weder, K.~Wild,
Quantum in the cloud: Application potentials and research opportunities,
in: Proc. 10th International Conference on Cloud Computing and Services Science
(CLOSER), SciTePress, 2020, pp. 9--24.

\bibitem{weder2021a}
B.~Weder, J.~Barzen, F.~Leymann, M.~Salm, D.~Vietz,
The quantum software lifecycle,
in: Proc. 1st ACM SIGSOFT International Workshop on Architectures and Paradigms
for Engineering Quantum Software (APEQS), ACM, 2020, pp. 2--9.

\bibitem{weder2021b}
% REPLACED: the entry previously here ("Quantum application archives:
% Supporting the reuse of quantum applications") could not be located in any
% bibliographic database and appears not to exist. Substituted with a real,
% topically adjacent paper by the same group -- VERIFY this is what you meant
% to cite, or delete this entry (it is currently uncited in the text).
B.~Weder, J.~Barzen, F.~Leymann, M.~Zimmermann,
Hybrid quantum applications need two orchestrations in superposition: A software
architecture perspective,
in: Proc. 18th IEEE International Conference on Web Services (ICWS), IEEE, 2021.

\bibitem{wild2021}
K.~Wild, U.~Breitenb\"{u}cher, L.~Harzenetter, F.~Leymann, D.~Vietz,
M.~Zimmermann,
TOSCA4QC: Two modeling styles for TOSCA to automate the deployment and
orchestration of quantum applications,
in: Proc. 24th IEEE International Enterprise Distributed Object Computing
Conference (EDOC), IEEE, 2020, pp. 125--134.

\bibitem{garcia2022}
J.~Garcia-Alonso, J.~Rojo, D.~Valencia, E.~Moguel, J.~Berrocal, J.~M. Murillo,
Quantum software as a service through a quantum API gateway,
IEEE Internet Computing 26~(1) (2022) 34--41.

\bibitem{beisel2023}
M.~Beisel, J.~Barzen, S.~Garhofer, F.~Leymann, F.~Truger, B.~Weder,
V.~Yussupov,
Quokka: A service ecosystem for workflow-based execution of variational quantum
algorithms,
in: Service-Oriented Computing -- ICSOC 2022 Workshops, Vol. 13821 of LNCS,
Springer, 2023, pp. 369--373.

\bibitem{quetschlich2023}
N.~Quetschlich, L.~Burgholzer, R.~Wille,
MQT Bench: Benchmarking software and design automation tools for quantum
computing,
Quantum 7 (2023) 1062.

\bibitem{li2020}
A.~Li, S.~Stein, S.~Krishnamoorthy, J.~Ang,
QASMBench: A low-level quantum benchmark suite for NISQ evaluation and
simulation,
ACM Transactions on Quantum Computing 4~(2) (2023) 1--26.

\bibitem{tomesh2022}
T.~Tomesh, P.~Gokhale, V.~Omole, G.~S. Ravi, K.~N. Smith, J.~Viszlai,
X.-C. Wu, N.~Hardavellas, M.~R. Martonosi, F.~T. Chong,
SupermarQ: A scalable quantum benchmark suite,
in: Proc. IEEE International Symposium on High-Performance Computer
Architecture (HPCA), IEEE, 2022, pp. 587--603.

\bibitem{finzgar2024}
J.~R. Fin\v{z}gar, P.~Ross, L.~H\"{o}lscher, J.~Klepsch, A.~Luckow,
QUARK: A framework for quantum computing application benchmarking,
in: Proc. IEEE International Conference on Quantum Computing and Engineering
(QCE), IEEE, 2022, pp. 226--237.

\bibitem{farhi2014}
E.~Farhi, J.~Goldstone, S.~Gutmann,
A quantum approximate optimization algorithm,
arXiv:1411.4028 (2014).

\bibitem{cerezo2021}
M.~Cerezo, A.~Arrasmith, R.~Babbush, S.~C. Benjamin, S.~Endo, K.~Fujii,
J.~R. McClean, K.~Mitarai, X.~Yuan, L.~Cincio, P.~J. Coles,
Variational quantum algorithms,
Nature Reviews Physics 3 (2021) 625--644.

\end{thebibliography}


\begin{thebibliography}{30}

\bibitem{nguyen2024}
H.~T. Nguyen, M.~Usman, R.~Buyya,
iQuantum: A toolkit for modeling and simulation of quantum computing
environments,
Software: Practice and Experience 54~(6) (2024) 1141--1171.

\bibitem{calheiros2011}
R.~N. Calheiros, R.~Ranjan, A.~Beloglazov, C.~A.~F. De~Rose, R.~Buyya,
CloudSim: A toolkit for modeling and simulation of cloud computing environments
and evaluation of resource provisioning algorithms,
Software: Practice and Experience 41~(1) (2011) 23--50.

\bibitem{mahmud2022}
R.~Mahmud, S.~Pallewatta, M.~Goudarzi, R.~Buyya,
iFogSim2: An extended iFogSim simulator for mobility, clustering, and
microservice management in edge and fog computing environments,
Journal of Systems and Software 190 (2022) 111351.

\bibitem{preskill2018}
J.~Preskill,
Quantum computing in the NISQ era and beyond,
Quantum 2 (2018) 79.

\bibitem{nielsen2010}
M.~A. Nielsen, I.~L. Chuang,
Quantum Computation and Quantum Information,
Cambridge University Press, Cambridge, UK, 2010.

\bibitem{cross2019}
A.~W. Cross, L.~S. Bishop, S.~Sheldon, P.~D. Nation, J.~M. Gambetta,
Validating quantum computers using randomized model circuits,
Physical Review A 100~(3) (2019).

\bibitem{proctor2025}
T.~Proctor, K.~Young, A.~D. Baczewski, R.~Blume-Kohout,
Benchmarking quantum computers,
Nature Reviews Physics 7 (2025) 105--118.

\bibitem{magesan2011}
E.~Magesan, J.~M. Gambetta, J.~Emerson,
Scalable and robust randomized benchmarking of quantum processes,
Physical Review Letters 106~(18) (2011).

\bibitem{magesan2012}
E.~Magesan, J.~M. Gambetta, B.~R. Johnson, C.~A. Ryan, J.~M. Chow,
S.~T. Merkel, M.~P. da~Silva, G.~A. Keefe, M.~B. Rothwell, T.~A. Ohki,
M.~B. Ketchen, J.~M. Steffen,
Efficient measurement of quantum gate error by interleaved randomized
benchmarking,
Physical Review Letters 109~(8) (2012).

\bibitem{yao2025}
G.~Yao, L.~Gupta,
A benchmarking framework for hybrid quantum--classical edge-cloud computing
systems,
Applied Sciences 15 (2025) 10245.

\bibitem{furutanpey2023}
A.~Furutanpey, J.~Barzen, M.~Bechtold, S.~Dustdar, F.~Leymann, P.~Raith,
F.~Truger,
Architectural vision for quantum computing in the edge-cloud continuum,
in: Proc. IEEE International Conference on Quantum Software (QSW), 2023,
pp. 88--103.

\bibitem{mastroianni2023}
C.~Mastroianni, L.~Scarcello, A.~Vinci,
Quantum computing management of a cloud/edge architecture,
in: Proc. ACM International Conference on Computing Frontiers (CF), 2023,
pp. 193--196.

\bibitem{xu2025}
M.~Xu, D.~Niyato, J.~Kang, Z.~Xiong, M.~Chen, D.~I. Kim, X.~Shen,
Hybrid reinforcement learning-based sustainable multi-user computation
offloading for mobile edge-quantum computing,
arXiv:2504.08134 (2025).

\bibitem{huang2024}
S.~Huang, L.~Wang, X.~Wang, B.~Tan, W.~Ni, K.-K. Wong,
Edge intelligence in satellite-terrestrial networks with hybrid quantum
computing,
arXiv:2409.19869 (2024).

\bibitem{satya2017}
M.~Satyanarayanan,
The emergence of edge computing,
Computer 50~(1) (2017) 30--39.

\bibitem{shi2016}
W.~Shi, J.~Cao, Q.~Zhang, Y.~Li, L.~Xu,
Edge computing: Vision and challenges,
IEEE Internet of Things Journal 3~(5) (2016) 637--646.

\bibitem{mao2017}
Y.~Mao, C.~You, J.~Zhang, K.~Huang, K.~B. Letaief,
A survey on mobile edge computing: The communication perspective,
IEEE Communications Surveys \& Tutorials 19~(4) (2017) 2322--2358.

\bibitem{abbas2018}
N.~Abbas, Y.~Zhang, A.~Taherkordi, T.~Skeie,
Mobile edge computing: A survey,
IEEE Internet of Things Journal 5~(1) (2018) 450--465.

\bibitem{leymann2020}
F.~Leymann, J.~Barzen, M.~Falkenthal, D.~Vietz, B.~Weder, K.~Wild,
Quantum in the cloud: Application potentials and research opportunities,
in: Proc. 10th International Conference on Cloud Computing and Services Science
(CLOSER), SciTePress, 2020, pp. 9--24.

\bibitem{weder2021a}
B.~Weder, J.~Barzen, F.~Leymann, M.~Salm, D.~Vietz,
The quantum software lifecycle,
in: Proc. 1st ACM SIGSOFT International Workshop on Architectures and Paradigms
for Engineering Quantum Software (APEQS), ACM, 2020, pp. 2--9.

\bibitem{weder2021b}
% REPLACED: the entry previously here ("Quantum application archives:
% Supporting the reuse of quantum applications") could not be located in any
% bibliographic database and appears not to exist. Substituted with a real,
% topically adjacent paper by the same group -- VERIFY this is what you meant
% to cite, or delete this entry (it is currently uncited in the text).
B.~Weder, J.~Barzen, F.~Leymann, M.~Zimmermann,
Hybrid quantum applications need two orchestrations in superposition: A software
architecture perspective,
in: Proc. 18th IEEE International Conference on Web Services (ICWS), IEEE, 2021.

\bibitem{wild2021}
K.~Wild, U.~Breitenb\"{u}cher, L.~Harzenetter, F.~Leymann, D.~Vietz,
M.~Zimmermann,
TOSCA4QC: Two modeling styles for TOSCA to automate the deployment and
orchestration of quantum applications,
in: Proc. 24th IEEE International Enterprise Distributed Object Computing
Conference (EDOC), IEEE, 2020, pp. 125--134.

\bibitem{garcia2022}
J.~Garcia-Alonso, J.~Rojo, D.~Valencia, E.~Moguel, J.~Berrocal, J.~M. Murillo,
Quantum software as a service through a quantum API gateway,
IEEE Internet Computing 26~(1) (2022) 34--41.

\bibitem{beisel2023}
M.~Beisel, J.~Barzen, S.~Garhofer, F.~Leymann, F.~Truger, B.~Weder,
V.~Yussupov,
Quokka: A service ecosystem for workflow-based execution of variational quantum
algorithms,
in: Service-Oriented Computing -- ICSOC 2022 Workshops, Vol. 13821 of LNCS,
Springer, 2023, pp. 369--373.

\bibitem{quetschlich2023}
N.~Quetschlich, L.~Burgholzer, R.~Wille,
MQT Bench: Benchmarking software and design automation tools for quantum
computing,
Quantum 7 (2023) 1062.

\bibitem{li2020}
A.~Li, S.~Stein, S.~Krishnamoorthy, J.~Ang,
QASMBench: A low-level quantum benchmark suite for NISQ evaluation and
simulation,
ACM Transactions on Quantum Computing 4~(2) (2023) 1--26.

\bibitem{tomesh2022}
T.~Tomesh, P.~Gokhale, V.~Omole, G.~S. Ravi, K.~N. Smith, J.~Viszlai,
X.-C. Wu, N.~Hardavellas, M.~R. Martonosi, F.~T. Chong,
SupermarQ: A scalable quantum benchmark suite,
in: Proc. IEEE International Symposium on High-Performance Computer
Architecture (HPCA), IEEE, 2022, pp. 587--603.

\bibitem{finzgar2024}
J.~R. Fin\v{z}gar, P.~Ross, L.~H\"{o}lscher, J.~Klepsch, A.~Luckow,
QUARK: A framework for quantum computing application benchmarking,
in: Proc. IEEE International Conference on Quantum Computing and Engineering
(QCE), IEEE, 2022, pp. 226--237.

\bibitem{farhi2014}
E.~Farhi, J.~Goldstone, S.~Gutmann,
A quantum approximate optimization algorithm,
arXiv:1411.4028 (2014).

\bibitem{cerezo2021}
M.~Cerezo, A.~Arrasmith, R.~Babbush, S.~C. Benjamin, S.~Endo, K.~Fujii,
J.~R. McClean, K.~Mitarai, X.~Yuan, L.~Cincio, P.~J. Coles,
Variational quantum algorithms,
Nature Reviews Physics 3 (2021) 625--644.

\end{thebibliography}
